\section{The strata algebra and the explicit class}\label{sec:strata}

This section fixes the relation space. The distinction between an
unnormalized gluing generator and an automorphism-weighted graph sum
will be used in both parts of the proof.

\subsection{Decorated graphs}

A stable graph \(\Gamma\) of type \((g,n)\) is a connected graph with
\(n\) labeled legs and genera \(g_v\geq0\) at its vertices, satisfying
\[
g=\sum_{v}g_v+|E(\Gamma)|-|V(\Gamma)|+1,
\qquad 2g_v-2+n_v>0.
\]
Here \(n_v\) counts all incident half-edges, including legs. Its gluing
map is
\[
\xi_\Gamma\colon\prod_{v\in V(\Gamma)}
\Mbar_{g_v,n_v}\longrightarrow\Mbar_{g,n}.
\]
At every half-edge let \(\psi\) denote the cotangent class. At a vertex
define
\[
\kappa_j=\pi_*(\psi_{n_v+1}^{j+1}),\qquad
\kappa_0=2g_v-2+n_v,
\]
where \(\pi\) forgets an additional marking.

The vector space \(\Salg_{g,n}\) is spanned by formal symbols
\([\Gamma,\gamma]\), modulo graph isomorphism, with \(\gamma\) a monomial
in the half-edge \(\psi\)'s and vertex \(\kappa_j\)'s for \(j\geq1\).
At each vertex its decoration degree is at most \(3g_v-3+n_v\).
The total degree is \(|E(\Gamma)|+\deg\gamma\), and
\[
q_{g,n}[\Gamma,\gamma]=\xi_{\Gamma*}(\gamma).
\]
There is no automorphism denominator in this definition of \(q_{g,n}\).
The per-vertex degree bound is the explicit convention in
\cite[Section~0.3]{PPZ}.

We use the usual decorated-graph intersection product. More explicitly,
take common refinements with specified contractions to fixed
representatives of the input graphs, requiring every refined edge to
survive in at least one contraction. Divide by the source symmetries
preserving those contractions. For each additional use of an overlapping
edge insert its normal factor \(-\psi-\psi'\); pull the leg and flag
decorations to their components and split kappa classes additively.
Specified contractions record the lifts to the gluing parameter spaces,
so this convention gives the product of the \emph{unnormalized}
pushforward generators. The construction is the transverse and
self-intersection rule for node cuts
\cite[Appendix~A, Proposition~9 and Equation~(11)]{GraberPandharipande2003}.

Terms exceeding any vertex dimension are set to zero. For multiplication
and restriction to a further graph splitting, an excessive decoration
remains excessive: the combined dimensions on the pieces decrease,
so at least one piece still exceeds its dimension.
Forgetting pullback of an undecorated graph is obtained by adding
the new leg successively at each vertex. These conventions lift the
corresponding Chow operations and do not impose any unproved relation.
The preceding dimension argument is not an assertion about forgetting,
which raises a vertex dimension. Compatibility of forgetting with the
formal detector is checked separately in Section~\ref{sec:detector}.

\subsection{The exact Pixton span}

Define
\[
A(T)=\sum_{m\geq0}\frac{(6m)!}{(3m)!(2m)!}T^m,\qquad
B(T)=\sum_{m\geq0}
\frac{6m+1}{6m-1}\frac{(6m)!}{(3m)!(2m)!}T^m.
\]
For a commuting parity variable \(\zeta^2=1\), put
\[
\widehat C_{3j}(T,\zeta)=T^jA(\zeta T),\qquad
\widehat C_{3j+1}(T,\zeta)=\zeta T^jB(\zeta T).
\]
Indices congruent to \(2\) modulo \(3\) are not used. For a series \(F\)
write
\[
\{F\}=\sum_{m\geq0,\ \epsilon\in\{0,1\}}
[F]_{T^m\zeta^\epsilon}K_{m,\epsilon}T^m.
\]
Give every vertex of \(\Gamma\) its own variable \(\zeta_v\), with
\(\zeta_v^2=1\). Define the operator
\[
\widehat\kappa_\Gamma
\left(\prod_{i=1}^{\ell}K_{e_i,\epsilon_i}\right)
=
\sum_{\tau\in\mathfrak S_\ell}
\prod_{\substack{c\text{ a cycle}\\\text{of }\tau}}
\left(
\sum_{v\in V(\Gamma)}
\kappa^{(v)}_{\sum_{i\in c}e_i}
\zeta_v^{\sum_{i\in c}\epsilon_i}
\right).
\]
The empty product maps to \(1\), and the operator acts coefficientwise
in \(T\). The permutation-cycle expression is Faber's formula, as
recorded in \cite[Equations~(1.12)--(1.13)]{ArbarelloCornalba1996}.

For an edge whose two half-edges lie at \(v,w\), with cotangent symbols
\(x,y\), set
\[
\Delta_e=
\frac{
A(\zeta_vxT)\zeta_wB(\zeta_wyT)
+\zeta_vB(\zeta_vxT)A(\zeta_wyT)
+\zeta_v+\zeta_w
}{(x+y)T}.
\]
This is a power series by
\[
A(T)B(-T)+A(-T)B(T)+2=0.
\]
The same rule applies to a loop.

\begin{definition}[Pixton's original span]\label{def:pixton-span}
Let \(\sigma=(\sigma_1,\ldots,\sigma_\ell)\) be a partition with no
part congruent to \(2\) modulo \(3\), possibly empty, and let
\(a_1,\ldots,a_n\geq0\) also avoid that residue class.
For \(0\leq r\leq3g-3+n\), require
\begin{equation}\label{eq:pixton-admissibility}
3r\geq g+1+|\sigma|+\sum_i a_i,\qquad
3r\equiv g+1+|\sigma|+\sum_i a_i\pmod2.
\end{equation}
Set
\begin{align}
P_\Gamma={}&
\left[
2^{-h^1(\Gamma)}
\widehat\kappa_\Gamma
\left(
\exp\{1-\widehat C_0\}
\prod_{j=1}^{\ell}\{\widehat C_{\sigma_j}\}
\right)\right.\notag\\
&\left.\hspace{12mm}\cdot
\prod_{i=1}^n\widehat C_{a_i}(\psi_iT,\zeta_{v(i)})
\prod_{e\in E(\Gamma)}\Delta_e
\right]_{T^{r-|E(\Gamma)|}\prod_v\zeta_v^{g_v+1}},
\label{eq:pixton-graph}
\end{align}
where \(h^1(\Gamma)=|E|-|V|+1\), \(v(i)\) carries the \(i\)-th leg,
and parity exponents are read modulo \(2\). In braces, the arguments
of \(\widehat C_i\) are \(T,\zeta\). A negative requested \(T\)-degree
means zero. The relation is
\[
\mathcal R(g,n,r;\sigma,a)=
\sum_\Gamma\frac{[\Gamma,P_\Gamma]}{|\Aut\Gamma|}.
\]

The space \(\Pspan_{g,n}\subseteq\Salg_{g,n}\) is the following
linear span in each degree. Choose an outer stable graph of type
\((g,n)\); put one such relation at one vertex, and arbitrary decorated
strata elements at all other vertices; then glue, using the
unnormalized gluing operation. Take the span over every choice,
including the graph with one vertex and no edges.
\end{definition}

\begin{proposition}[Known validity and stability]\label{prop:known-relations}
The span in Definition~\ref{def:pixton-span} is an ideal and is stable
under tautological gluing and forgetting pushforwards and pullbacks.
For gluing pullback to a product, the target relation space is the sum
of the spaces with a relation in one factor.
Moreover
\[
\Pspan_{g,n}\subseteq\ker q_{g,n}.
\]
\end{proposition}

\begin{proof}
The ideal and operation statements are
\cite[Proposition~1]{Pixton} for precisely this original span.
Chow validity of the \(a_i\in\{0,1\}\) seed relations follows from
\cite[Theorem~1 and its reconstruction]{Janda}.
The construction in \cite[Section~3.5]{PPZ} generates the higher
marking indices by psi multiplication, the partition insertions by
forgetting extra markings, and the outer-graph relations by gluing.
These are formal identities, so they apply to the Chow-valid seeds.
The parity condition in \eqref{eq:pixton-admissibility} converts the
strict seed inequality to the displayed admissibility inequality.
In Janda's reconstructed normalization, a degree-\(r\) seed is
\((-1)^{\sum_i a_i}(-1/1728)^r\) times the empty-partition recipe
above. This nonzero scalar does not change its vanishing or its span.
\end{proof}

No converse to Proposition~\ref{prop:known-relations} is used.
In particular, same-vertex multiplication of a relation does not enlarge
the span, whereas an arbitrary new Chow relation is not assumed to
belong to it.

\subsection{The corrected cubic and its alternation}

For the rest of the construction fix
\begin{equation}\label{eq:global-parameters}
g=10^{60},\qquad D_*=\binom g3,\qquad n=3D_*.
\end{equation}
Form a word \(\mathbf i\) in the colors \(1,\ldots,g\) by concatenating
all increasing triples of colors in lexicographic order. Its consecutive
three-slot blocks will be denoted \(J\). Each color occurs
\(\binom{g-1}{2}\) times.

For a set of labels \(T\) of size at least two, let \(D_T\) be the
one-edge class cutting off a genus-zero tail with precisely those
labels. For distinct labels \(a,b\), put
\[
\bar D_{ab}=\sum_{T\supseteq\{a,b\}}D_T.
\]
We also need a class supported on rational bridges, not tails.

\begin{definition}[The explicit counterexample]\label{def:candidate}
On three labeled slots \(J=(a,b,c)\), take every one-edge
\emph{unmarked} stable graph of genus \(g\), subdivide its edge by
a genus-zero vertex, and put all three labels at that new vertex.
Sum the resulting undecorated graphs, dividing each by its automorphism
order. Pull this sum back by forgetting all labels outside \(J\);
call the resulting class \(\beta_J^{\mathrm{st}}\).

Define
\begin{equation}\label{eq:candidate}
F_J^{\mathrm{st}}=
\bar D_{ab}\bar D_{ac}-\frac12\beta_J^{\mathrm{st}},
\qquad
Y=\sum_{\substack{\tau\in\mathfrak S_n\\\tau\mathbf i=\mathbf i}}
\operatorname{sgn}(\tau)\,
\tau\left(\prod_J F_J^{\mathrm{st}}\right).
\end{equation}
Here \(\tau\) relabels slots, and products and forgetting pullbacks
are the formal graph operations fixed above.
\end{definition}

Every operation in Definition~\ref{def:candidate} is finite.
Each cubic factor has codimension two, so
\[
\deg Y=2D_*<3g-3+n.
\]
The pair \((g,n)\) is stable. The automorphism weights in
\(\beta_J^{\mathrm{st}}\) specify coefficients of this particular
class; they do not change the unnormalized generators of
\(\Salg_{g,n}\).

The proof of Theorem~\ref{thm:main} now divides into two independent
claims. Proposition~\ref{prop:detector} proves formal nonmembership.
The remaining sections establish the geometric vanishing
in Theorem~\ref{thm:chow-vanishing}.
